\documentclass[11pt]{article}

\usepackage[utf8]{inputenc}
\usepackage[T1]{fontenc}
\usepackage[spanish,es-tabla]{babel}
\usepackage{lmodern}
\usepackage[a4paper,margin=2cm]{geometry}
\usepackage{graphicx}
\usepackage{float}
\usepackage{xcolor}
\usepackage{colortbl}
\usepackage{array}
\usepackage{enumitem}
\usepackage{longtable}
\usepackage{booktabs}
\usepackage{hyperref}
\hypersetup{colorlinks=true,linkcolor=black,urlcolor=[HTML]{0B3A5C}}

\graphicspath{{capturas/}}

% Enlaces al prototipo publicado y al código en GitHub.
\newcommand{\urlapp}{https://thankstochatgpt.github.io/LOARD}
\newcommand{\app}[1]{\href{\urlapp#1}{\texttt{#1}}}
\newcommand{\codigo}[1]{\href{https://github.com/ThanksToChatGPT/LOARD/tree/main/Hito\%202/Modulo-I-Prototipo/prototipo/src/#1}{\texttt{src/#1}}}
\setlength{\parskip}{0.5em}
\setlength{\parindent}{0pt}
\renewcommand{\arraystretch}{1.25}
\setlist{nosep,leftmargin=1.2em}

\definecolor{marca}{HTML}{0B3A5C}
\definecolor{borde}{HTML}{B8C4CE}
\arrayrulecolor{borde}
\newcommand{\prop}[1]{\texttt{#1}}
\newcommand{\imagen}[2][\linewidth]{%
  \fcolorbox{borde}{white}{%
    \includegraphics[width=\dimexpr#1-2\fboxsep-2\fboxrule\relax,height=9.5cm,keepaspectratio]{#2.png}%
  }%
}
% Contenido de la ficha: qué es, variantes y estados.
\newcommand{\ficha}[3]{%
  {\large\bfseries\color{marca}#1}\par\smallskip
  \textbf{Qué es.} #2\par\smallskip
  #3}

% Ficha con la imagen a la izquierda (55% imagen, 45% texto).
% #1 imagen, #2 nombre, #3 qué es, #4 variantes y estados
\newcommand{\comp}[4]{%
  \par\medskip\noindent
  \begin{minipage}{\linewidth}
  \begin{tabular}{|m{0.55\linewidth}|m{\dimexpr0.45\linewidth-4\tabcolsep-3\arrayrulewidth\relax}|}
  \hline
  \vspace{4pt}\centering\imagen{#1}\vspace{4pt} & \vspace{4pt}\ficha{#2}{#3}{#4}\vspace{4pt} \\
  \hline
  \end{tabular}
  \end{minipage}\par\medskip}

% Ficha con la imagen arriba (componentes anchos).
% #1 imagen, #2 nombre, #3 qué es, #4 variantes y estados
\newcommand{\compancho}[4]{%
  \par\medskip\noindent
  \begin{minipage}{\linewidth}
  \begin{tabular}{|m{\dimexpr\linewidth-2\tabcolsep-2\arrayrulewidth\relax}|}
  \hline
  \vspace{4pt}\centering\imagen{#1}\vspace{4pt} \tabularnewline
  \hline
  \vspace{4pt}\ficha{#2}{#3}{#4}\vspace{4pt} \\
  \hline
  \end{tabular}
  \end{minipage}\par\medskip}

\title{Inventario de componentes\\[0.3em]
\large Loard - Hito 2 - Módulo II: Arquitectura de frontend y componentización}
\author{}
\date{}

\begin{document}
\maketitle

\section{Introducción}

Para el desarrollo de los prototipos, se plantearon componentes que se repetirían varias veces, como botones, formularios, etiquetas, tablas y avisos. Para mantener consistencia en todo el prototipo y agilizar el trabajo en equipo, se inició creando los componentes \textbf{una sola vez} para luego reutilizarlos en todos los módulos sin crear elementos visuales propios.

Este inventario muestra cada componente con una captura tomada del prototipo en funcionamiento. El código se encuentra en \codigo{components} y se recomienda ver la página de muestra del prototipo para visualizar estos componentes
en la ruta \app{/muestra} del prototipo publicado.

La sección~\ref{sec:react} explica cómo se organizan en React + TypeScript (carpetas, jerarquía y props) y la sección~\ref{sec:historias} muestra qué componentes usa cada historia de usuario.


\subsection*{Reglas comunes a todos los componentes}

\begin{itemize}
  \item \textbf{Tamaño} (\prop{size}): \prop{default} para escritorio y \prop{field} para uso en obra (56 px de alto en móvil, más fácil de tocar).
  \item \textbf{Tono} (\prop{tone}): un mismo juego de colores con significado fijo: \prop{info} (azul), \prop{success} (verde), \prop{warning} (ámbar), \prop{error} (rojo) y \prop{neutral}.
  \item \textbf{Tema}: todos funcionan en claro y oscuro sin cambios.
  \item \textbf{Bloqueado no es deshabilitado}: un control bloqueado sigue visible y \emph{explica el motivo} del bloqueo.
\end{itemize}

\newpage
\section{Átomos}

\comp{button}{Button }
{El botón de acción. Es el componente más repetido del sistema.}
{\textbf{Variantes:} primario (acción principal), secundario, advertencia (acciones delicadas como ``Devolver'') y fantasma (acciones de poco peso). 
\textbf{Estados:} normal, con el cursor encima, con foco, cargando (muestra un spinner y cambia el texto) y bloqueado (gris, con el motivo debajo).}

\comp{badges}{Etiquetas de estado}
{Etiquetas pequeñas con ícono, color y texto. \textbf{StatusBadge} muestra un estado oficial del negocio; \textbf{Badge} es la etiqueta genérica para cualquier otro dato corto.}
{\textbf{Variantes de StatusBadge:} una por cada estado oficial (estados SST del personal, estados del corte y estados de los requisitos). \textbf{Tonos de Badge:} info, éxito, advertencia, error y neutro. El estado nunca se comunica solo con color: siempre lleva ícono y texto.}

\comp{checkbox-chip}{Checkbox y Chip}
{\textbf{Checkbox}: casilla para marcar una opción (por ejemplo, seleccionar personal del corte). \textbf{Chip}: pastilla que se oprime para elegir o insertar un texto rápido.}
{\textbf{Checkbox:} sin marcar, marcada y bloqueada (no se deja marcar y la fila se sacude para avisar). \textbf{Chip:} normal y seleccionado.}

\comp{progreso}{ProgressBar, Spinner, Skeleton e IconButton}
{\textbf{ProgressBar}: barra de avance. \textbf{Spinner}: indicador giratorio de 16 px, solo dentro de botones. \textbf{Skeleton}: franja gris que ocupa el lugar del contenido mientras carga. \textbf{IconButton}: botón cuadrado de 48 x 48 px con solo un ícono.}
{\textbf{ProgressBar:} con valor, con tono de advertencia e indeterminada (cuando no se conoce el avance). \textbf{IconButton:} con o sin contador.}

\newpage
\section{Moléculas}

\comp{campos}{TextField, DateField y TextArea (campos de formulario)}
{Campos para escribir. Todos comparten la misma estructura: etiqueta arriba, el campo, y debajo una ayuda o el mensaje de error. \textbf{DateField} es para fechas y \textbf{TextArea} para textos largos.}
{\textbf{Estados:} vacío, con foco, lleno, obligatorio (marcado con *) y con error (borde rojo y mensaje que dice qué pasó y cómo corregirlo).}

\comp{seleccion}{Select, Combobox, SearchField, RadioCardGroup y FilterChips (selección)}
{Formas de elegir entre opciones. \textbf{Select}: lista desplegable corta. \textbf{Combobox}: lista en la que se puede escribir para buscar. \textbf{SearchField}: buscador. \textbf{RadioCardGroup}: opciones excluyentes en tarjetas grandes. \textbf{FilterChips}: filtros rápidos con contador.}
{\textbf{Estados:} sin selección, con opción elegida y desplegado (Select y Combobox). En FilterChips, el filtro activo queda resaltado.}

\compancho{kpi}{KpiCard (tarjeta de indicador)}
{Tarjeta que muestra una cifra importante con su título y una descripción. Se usa en el tablero.}
{\textbf{Variantes:} informativa, con cifra en color de estado y como filtro (se puede oprimir). \textbf{Estados:} normal y seleccionada.}

\comp{banner}{Banner}
{Franja de aviso dentro de la pantalla. Informa algo que el usuario debe saber antes de continuar.}
{\textbf{Variantes por tono:} información, advertencia, error o bloqueo, y confirmación. Cada una con su color e ícono.}

\comp{conexion-upload}{ConnectionIndicator y UploadItem}
{\textbf{ConnectionIndicator}: muestra en la barra superior si hay señal. \textbf{UploadItem}: una fila por cada archivo que se está subiendo.}
{\textbf{ConnectionIndicator:} en línea, sin señal, sincronizando y error. \textbf{UploadItem:} en espera, subiendo (con barra), confirmando, cargado y con error (permite reintentar).}

\comp{toast}{Aviso temporal}
{Mensaje breve que aparece unos segundos en la parte inferior para confirmar una acción.}
{\textbf{Variantes:} informativo y con opción ``Deshacer''. Se cierra solo o con la X.}

\newpage
\section{Organismos}

\compancho{tabla}{ResponsiveTable}
{La tabla de todas las bandejas y listados. En escritorio se ve como tabla; en móvil cada fila se convierte en una tarjeta, sin desplazamiento horizontal.}
{\textbf{Variantes:} tabla (escritorio) y tarjetas (móvil). \textbf{Estados:} con datos, primera carga (filas Skeleton), recarga en segundo plano y fila resaltada en rojo.}

\comp{tabla-movil}{ResponsiveTable en móvil}
{La misma tabla anterior vista en un teléfono. El componente y los datos son los mismos; solo cambia la presentación.}
{\textbf{Presentación:} cada fila es una tarjeta con sus datos en vertical.}

\comp{sidebar}{Sidebar}
{Menú de navegación principal, fijo a la izquierda en escritorio.}
{\textbf{Estados de cada opción:} normal, con el cursor encima y activa (resaltada con una barra a la izquierda).}

\compancho{topbar}{Topbar}
{Barra fija arriba con el proyecto actual, el estado de conexión, el rol del usuario y accesos rápidos.}
{\textbf{Estados:} cambia según el estado de conexión y la cantidad de alertas pendientes.}

\compancho{offline-banner}{Franja sin conexión}
{Franja bajo la barra superior que avisa cuando se pierde la señal en obra y cuántos cambios quedan guardados en el dispositivo.}
{\textbf{Estados:} sin conexión, sincronizando, todo enviado (se oculta sola a los 5 segundos) y envío rechazado (pide revisar).}

\comp{dialog}{Diálogo de confirmación}
{Ventana que aparece sobre la pantalla para confirmar una acción importante o pedir un dato (por ejemplo, la observación al devolver un corte).}
{\textbf{Variantes:} confirmación simple y con formulario. \textbf{Estados:} cerrada, abierta y ``¿descartar cambios?''.}

\compancho{drawer-full}{Panel lateral}
{Panel que se desliza desde la derecha para ver un detalle o un documento sin salir de la pantalla. En móvil ocupa toda la pantalla.}
{\textbf{Variantes:} panel lateral (desde 1024 px) y pantalla completa. \textbf{Estados:} cerrado y abierto.}

\comp{fileupload}{FileUpload}
{Zona para arrastrar o seleccionar soportes. Debajo lista cada archivo con un UploadItem.}
{\textbf{Estados:} vacía, con un archivo encima (resaltada) y con archivos listados.}

\comp{timeline}{Timeline}
{Línea de tiempo con el historial de cambios de un registro: quién, cuándo y de qué estado a cuál.}
{\textbf{Detalle:} cada registro puede llevar o no un motivo.}

\newpage
\section{Plantillas}

\compancho{appshell}{Marco de página}
{\textbf{AppShell} es el marco común de todas las pantallas: menú lateral, barra superior y zona de contenido. \textbf{PageHeader} es el encabezado de cada pantalla: título, descripción y acciones.}
{\textbf{Variantes:} escritorio (menú lateral fijo) y móvil (navegación compacta). Con o sin franja de aviso.}

\comp{formlayout}{Plantilla de formulario}
{Estructura común de todos los formularios: encabezado, campos en una columna centrada y un pie fijo con los botones.}
{\textbf{Presentación:} en escritorio ocupa 8 de 12 columnas centradas; el pie con los botones queda siempre visible.}

% =====================================================================
\newpage
\section{Estructura en React + TypeScript}\label{sec:react}

\subsection*{Organización de las carpetas}

El prototipo es una aplicación de React 19 con TypeScript, Vite y Tailwind CSS. Los componentes siguen el diseño atómico: átomos, moléculas, organismos y plantillas. Todos se importan desde un solo punto, \texttt{@/components}.

\begin{center}
\begin{tabular}{|>{\raggedright\arraybackslash}m{3.6cm}|>{\raggedright\arraybackslash}m{12cm}|}
\hline
\textbf{Carpeta} & \textbf{Qué va ahí} \\ \hline
\codigo{styles} & \texttt{tokens.css} es el único archivo con códigos de color. En el resto se usan nombres de token (\texttt{bg-surface}, \texttt{text-error}). \\ \hline
\codigo{components} & Componentes compartidos por todo el equipo, en \texttt{atoms}, \texttt{molecules}, \texttt{organisms} y \texttt{templates}. \\ \hline
\codigo{lib} & Reglas y constantes de negocio: estados oficiales, semáforo SST, formatos de fecha, validación de archivos, datos SST compartidos y consultas de seguimiento. \\ \hline
\codigo{hooks} & Tema claro u oscuro, foco dentro de diálogos, ancho de pantalla y estado de conexión. \\ \hline
\codigo{features} & Una carpeta por módulo funcional: \texttt{acceso}, \texttt{cortes}, \texttt{dossier}, \texttt{sst} y \texttt{tablero}. Cada pantalla se construye dentro de la suya. \\ \hline
\codigo{muestra} & Página que enseña todos los componentes compartidos (\app{/muestra}). \\ \hline
\end{tabular}
\end{center}


\subsection*{Props: el ejemplo de \texttt{Button}}

Cada componente declara sus props con una interfaz de TypeScript. Las variantes y los estados son props, no clases sueltas. Así es la de \texttt{Button}:

\begin{center}
\begin{tabular}{|>{\ttfamily}m{3cm}|>{\ttfamily\raggedright\arraybackslash}m{4.4cm}|>{\raggedright\arraybackslash}m{7.8cm}|}
\hline
\textrm{\textbf{Prop}} & \textrm{\textbf{Tipo}} & \textbf{Para qué sirve} \\ \hline
variant & 'primary' | 'secondary' | 'warning' | 'ghost' & Peso visual de la acción. \\ \hline
size & 'default' | 'field' & \texttt{field} sube el alto a 56 px en móvil para uso en obra. \\ \hline
icon & ReactNode & Ícono a la izquierda del texto. \\ \hline
loading & boolean & Petición en curso: muestra el spinner e impide un segundo envío. \\ \hline
loadingText & string & Texto mientras carga, por ejemplo ``Guardando corte\ldots{}''. \\ \hline
blocked & boolean & Bloqueado por una regla de negocio; sigue recibiendo foco. \\ \hline
blockedReason & ReactNode & Motivo visible del bloqueo y cómo resolverlo. Obligatorio si \texttt{blocked}. \\ \hline
fullWidth & boolean & Ocupa todo el ancho (botones del pie en móvil). \\ \hline
\end{tabular}
\end{center}

Los demás componentes siguen el mismo patrón: \texttt{Banner} recibe \texttt{tone}, \texttt{title} y \texttt{actions}; \texttt{Dialog} y \texttt{Drawer} reciben \texttt{open}, \texttt{onClose} y \texttt{title}; \texttt{PageHeader} recibe \texttt{title}, \texttt{description}, \texttt{meta} y \texttt{actions}; \texttt{FormLayout} recibe \texttt{header}, \texttt{status}, \texttt{footer} y \texttt{onSubmit}.

\subsection*{Jerarquía de una pantalla: el detalle del corte}

Las pantallas componen las piezas compartidas y las reglas viven fuera de los componentes. Por ejemplo, el detalle del corte (\texttt{/cortes/:id}) se arma así:

\begin{verbatim}
DetalleCorte (pagina /cortes/:id)
|-- PageHeader . StatusBadge . AvisoDevolucion
|-- ValoresCorte (actividades, horas-hombre, insumos)
|-- SoportesCorte
|   |-- ItemSoporte x n  -->  VisorSoporte (Drawer)
|   `-- FileUpload  -->  UploadItem x n
`-- RevisionCorte (solo Director, corte en revision)
    |-- Dialog "Confirmar visado tecnico"
    `-- DialogoObservacion  -->  TextArea + Chip x 5
\end{verbatim}

Las reglas están en \texttt{soportes.ts} (validación, permisos y retiro sin borrado) y \texttt{modeloDemo.ts} (quién puede decidir, sobre qué versión y con qué observación). \texttt{CortesDemoContext} las comparte con las pantallas. Del mismo modo, en el dossier las reglas están en \texttt{modelo.ts} y \texttt{entrega.ts}, y \texttt{lib/seguimiento} lleva esos datos a la bitácora y al tablero sin crear otra checklist.

% =====================================================================

Con estos 33 componentes se construyen las pantallas del MVP. Si cambia una regla de diseño (un color, un tamaño), se corrige en un solo componente y el cambio llega a todas las pantallas.

\end{document}
